% generated by roam.py handbook export -- не редактировать, править заметки
% 20261013224457-квадратичная_сходимость_сверхлинеина.org (body cbc7622ce124)

\begin{lemma}[{Квадратичная сходимость сверхлинейна}]\label[lemma]{note:4CD0DF77-0553-4E99-901D-4D4C50C9900C}
  Пусть \hyperref[note:C42D2C19-B35F-4C0E-AFEB-D880CBB72176]{невязка}  это $r_k$.
  Если $r_{k+1} \le C r_k^2$,  (\hyperref[note:FDA5DA6F-DE0F-437E-A4F4-89757BA7F538]{Квадратичная сходимость}) и $Cr_0 < 1$ \\
  То сходимость~--- \hyperref[note:9FA20DA5-3A95-4683-9CB5-C6C9A903C86A]{сверхлинейная}
\end{lemma}

\begin{proof}
  Воспользуемся \hyperref[note:C0408924-F658-43DF-8794-0C9AA85E7580]{тестом корней} в форме отношений: достаточно $\frac{r_{k+1}}{r_k} \rightarrow 0$.
  Заметим, что $\frac{r_{k+1}}{r_k}\le \frac{C r^2_k}{r_k} = C r_k$ \\
  Тогда для $\frac{r_{k+1}}{r_k} \rightarrow 0$ достаточно доказать, что $Cr_{k} \rightarrow 0$.

  Теперь давайте докажем, что $r_k \rightarrow 0$.

  Докажем убывание $r_k$ по мат индукции.
  Пусть $q:=Cr_0$, $q < 1$ по условиям.

  Доказываем $r_k \le q^k r_0$. База $r_1 \le Cr_0^2 = (Cr_0) r_0 = q r_0$

  Переход
  \[
  r_{k+1} \le Cr_k^2 = C r_k r_k \le C r_k q^k r_0 = \underbrace{(C r_0)}_{q} q^k r_k = q^{k+1} r_k
  \]
  Далее заметим, что $r_k \le q^k r_0$, а $q^k < 1$, тогда $r_k \le q^k r_0\le r_0$ \\
  Давайте завершим нашу оценку
  \[
  r_{k+1} \le q^{k+1}r_k \le q^{k+1}r_0
  \]
  Доказано.

  Тогда
  \[
  0 \le \frac{r_{k+1}}{r_k}\le Cr_k \le Cq^kr_0 \rightarrow 0
  \]
  А это и есть сверхлинейная сходимость;
  ЧТД
\end{proof}
